% =============================================================================
% AGENT   : Statistician (outputs) + Geneticist (interpretation)
% MODULE  : Results — English
% =============================================================================

\section{Results}
\label{sec:results}

\subsection{Overview of Simulated Marker Panel}

A total of 60 molecular markers were evaluated, comprising 36 \snp{} and
24 \ssr{} markers. Descriptive statistics for \pic{}, \he{}, and \ho{}
are presented in \cref{tab:marker_summary} (Supplementary Material).

\subsection{Polymorphism Information Content}

\snp{} markers displayed a mean \pic{} of approximately 0.25--0.32,
reflecting the inherent constraint of biallelic variation. In contrast,
\ssr{} markers modelled with symmetric allele frequency distributions
exhibited substantially higher mean \pic{} values (0.35--0.42), consistent
with their established role as highly polymorphic marker systems
\autocite{Powell1996}.

The distribution of \pic{} values is illustrated in
\cref{fig:pic_boxplot}. \ssr{} markers showed a wider interquartile
range, indicating greater variability in informativeness across loci.

\subsection{Heterozygosity}

Expected heterozygosity (\he{}) followed the same pattern as \pic{},
as both statistics are direct functions of allele frequencies. The
mean \he{} for \ssr{} markers exceeded that of \snp{} markers,
reflecting the higher balanced allele frequencies in the former.
Observed heterozygosity (\ho{}) values were closely distributed around
\he{} (mean deviation $\leq 0.04$), suggesting no major departure from
Hardy--Weinberg equilibrium in the simulated population.
